%\pagebreak
\section{\kcode{split} Construct}
\label{sec:split}
\index{constructs!split@\kcode{split}}
\index{split construct@\kcode{split} construct}

%\index{counts!full@\kcode{counts}}
\index{counts clause@\kcode{counts} clause}
\index{clauses!counts@\kcode{counts}}

\index{OpenMP identifiers!omp_fill@\kcode{omp_fill}}
\index{omp_fill identifier@\kcode{omp_fill} identifier}

The \kcode{split} construct splits a single canonical loop having an iteration
count of \plc{n} into
multiple loops by partitioning the iteration space among \plc{m} loops
where \plc{m} is the number of arguments in the \kcode{counts} clause.
For each argument of the \kcode{counts} clause a loop is 
generated with the number of iterations specified by the argument value.
The original iteration space is sequentially partitioned among the loops.

The special \kcode{omp_fill} identifier specifies
the remaining number of iterations after all other iterations have been accounted for.
That is, the \kcode{omp_fill} value is equal to the total number of iterations less the
sum of all the other iteration counts.

Examples in this section contain \emph{equivalent}
routines that illustrate the functionality of the constructs. 
The equivalent routines may undergo further compiler optimization.

In the following example the \kcode{split} construct splits the loop
in the \ucode{adder} routine
into 3 loops of iteration counts \ucode{m}, \ucode{1}, and
\ucode{n-m-1}, to allow loop iterations to be vectorized in the first and last loops.

    \cexample[6.0]{split}{1}
\ffreeexample[6.0]{split}{1}
\clearpage

%%%%%%%%%%%%%%%%%%%%%%%%%%%%%%%%%%%%%%%%%%%%%%%%%%%%%%%%%%%%%%%%%%%%%%%%
The next example illustrates the concerted application
of two loop-transforming constructs (\kcode{split} and \kcode{fuse})
on a single loop. The \kcode{split} construct splits the iteration
space among three loops with the \kcode{counts(omp_fill,\ucode{n},\ucode{n})} clause,
with iterations counts of \ucode{1}, \ucode{n} and \ucode{n}. This effectively removes the
dependence introduced by the \ucode{if(...)} conditional, and generates
two loops spanning the next \ucode{n} iterations and the following \ucode{n} iterations.
The \kcode{fuse} construct combines the latter two
loops, so that the two element-updates in each iteration are to the same
position in the array. The fusion increases memory locality.

    \cexample[6.0]{split}{2}
\ffreeexample[6.0]{split}{2}

%%%%%%%%%%%%%%%%%%%%%%%%%%%%%%%%%%%%%%%%%%%%%%%%%%%%%%%%%%%%%%%%%%%%%%%%
% The next example illustrates splitting a loop between a host and a device,
% with the assistance of the \kcode{apply} clause.
% A loop with an iteration count \ucode{n} is split into two loops of \ucode{m} and 
% \ucode{n-m} counts by the \kcode{counts(\ucode{m},omp_fill)} clause. The
% \kcode{apply} clause specifies that the \kcode{target}~\kcode{loop} offloading
% directive is applied to the first loop, and the \kcode{parallel do|for}
% directive is applied to the second loop for execution on the host.
% 
% %The \ucode{split_composable_equivalent} routine shows a simple 
% %semantic equivalent of the \kcode{split} construct implementation.
% 
%     \cexample[6.0]{split}{3}
% \ffreeexample[6.0]{split}{3}

% %%%%%%%%%%%%%%%%%%%%%%%%%%%%%%%%%%%%%%%%%%%%%%%%%%%%%%%%%%%%%%%%%%%%%%%%
% The final example of this section illustrates splitting a loop to
% run across different offload devices.
% The \kcode{counts(n/2,omp_fill)} clause splits the n-iteration loop
% into two loops having \kcode{n/2} and \kcode{n-n/2} iterations.
% The \kcode{apply} clause applies the \kcode{target loop}
% directive, and specifies
% two different devices, \ucode{device(0)} and \ucode{devices(1)}
% for execution of the two loops.
% 
% The \ucode{split_composable_equivalent} routine shows a simple 
% semantic equivalent of the \kcode{split} construct implementation.
% 
% !!! NEEDS a wait (is auto mapping ok?)
% 
%     \cexample[6.0]{split}{4}
% \ffreeexample[6.0]{split}{4}
